% ============================================================
%  CHAPTER 17: WAVE MOTION
%  The Physics Codex: A Complete University Foundation
%  Korede Samuel Omotosho (Falcon)
%  Aurikrex Learning Systems, 2026
%
%  File: chapters/part3/ch17_wave_motion.tex
% ============================================================

\chapter{Wave Motion}
\label{ch:wave_motion}
\minitoc
\newpage

\chapterquote{The career of a young theoretical physicist
consists of treating the harmonic oscillator in ever
increasing levels of abstraction.}{Sidney Coleman}

\begin{objectivebox}
By the end of this chapter, you should be able to:
\begin{enumerate}[noitemsep]
  \item Define a wave and explain what is transferred
  in wave motion
  \item Distinguish between transverse and
  longitudinal waves with examples
  \item Define amplitude, wavelength, frequency,
  period, wave speed and phase
  \item Apply the wave equation $v = f\lambda$
  \item Describe the properties of waves: reflection,
  refraction, diffraction and interference
  \item State the principle of superposition and
  apply it to interference
  \item Describe stationary (standing) waves and
  explain how they form
  \item Calculate the positions of nodes and
  antinodes in standing waves
  \item Explain resonance in strings and pipes
\end{enumerate}
\end{objectivebox}

% ============================================================
\section{Intuition: What is a Wave?}
% ============================================================

Drop a stone into a still pond. Circular ripples
spread outward from the point of impact. Watch a
cork floating on the surface: as the ripples pass,
the cork bobs up and down but does not travel outward
with the ripples. The water molecules move up and
down; they do not travel horizontally. Yet the
pattern of the ripple --- the disturbance --- travels
outward across the entire pond.

This is the essence of wave motion. A wave is a
\textbf{travelling disturbance} that transfers
\textbf{energy} from one place to another
\textit{without transferring matter}. The medium
through which the wave travels oscillates about its
equilibrium position, but the medium itself does
not travel with the wave.

Waves are everywhere. Sound, light, ocean tides,
radio signals, earthquakes, the vibrations of a
guitar string, the quantum mechanical behaviour of
electrons --- all are described by wave mathematics.
Chapter~17 establishes the language of waves that
the rest of Part 3 will use.

% ============================================================
\section{Types of Waves}
% ============================================================

\begin{definitionbox}[Transverse and Longitudinal Waves]
\textbf{Transverse wave:} The oscillation of the
medium is \textit{perpendicular} to the direction
of wave travel (energy propagation).

Examples: light, electromagnetic waves, waves on
a rope or string, water surface waves, S-waves
(seismic shear waves).

\vspace{0.4cm}

\textbf{Longitudinal wave:} The oscillation of the
medium is \textit{parallel} to (along) the direction
of wave travel.

Examples: sound waves in air, P-waves (seismic
pressure waves), compression waves in a spring.

\vspace{0.4cm}

\textbf{Key distinction:} Transverse waves can be
polarised (oscillations restricted to one plane);
longitudinal waves cannot be polarised.
\end{definitionbox}

\begin{diagbox}[Transverse vs Longitudinal Waves]
\begin{center}
\begin{tikzpicture}[scale=1.0, font=\small]

  % === TRANSVERSE WAVE ===
  \begin{scope}[shift={(0,5.5)}]
    \node[font=\bfseries] at (6,2.5)
      {Transverse Wave};
    \node[font=\small, align=center] at (6,2.0)
      {(oscillation $\perp$ direction of travel)};

    % Wave direction arrow
    \draw[->, very thick, codexgold]
      (0.5,1.7) -- (2.5,1.7)
      node[right, font=\scriptsize]{Direction of energy travel};

    % Transverse wave shape
    \draw[codexblue, very thick, domain=0:12,
      samples=100]
      plot (\x*0.95, {0.8*sin(deg(\x*0.52))+0.5});

    % Oscillation arrows (perpendicular)
    \foreach \x in {0,1,2,...,11} {
      \pgfmathsetmacro\yval{0.8*sin(deg(\x*0.52*1.05))+0.5}
      \pgfmathsetmacro\sign{sign(cos(deg(\x*0.52*1.05)))}
      \draw[->, thick, codexred!50, line width=0.8pt]
        ({\x*0.95},0.5)
        -- ({\x*0.95}, \yval);
    }

    % Labels
    \draw[<->, gray] (0,1.3) -- (0,-0.3);
    \node[left, gray, font=\tiny, align=center] at (0,0.5)
      {Oscillation\\direction};
  \end{scope}

  % === LONGITUDINAL WAVE ===
  \begin{scope}[shift={(0,0)}]
    \node[font=\bfseries] at (6,3.5)
      {Longitudinal Wave};
    \node[font=\small, align=center] at (6,3.0)
      {(oscillation $\parallel$ direction of travel)};

    % Wave direction arrow
    \draw[->, very thick, codexgold]
      (0.5,1.5) -- (2.5,1.5)
      node[right]{Direction of energy travel};

    % Particles arranged horizontally
    \foreach \i in {0,...,23} {
      % Compression/rarefaction pattern
      \pgfmathsetmacro\phase{\i*15}
      % pgfmath's sin() uses degrees by default; deg() here would blow up the coordinate.
      \pgfmathsetmacro\offset{0.35*sin(\phase)}
      \pgfmathsetmacro\xpos{\i*0.50 + \offset}
      \fill[codexblue] (\xpos, 2.2) circle (4pt);
    }

    % Compression label
    \draw[<->, codexred, thick]
      (1.6,0.8) -- (3.0,0.8);
    \node[below, codexred, font=\tiny, align=center]
      at (2.3,0.7) {Compression\\(high density)};

    % Rarefaction label
    \draw[<->, codexblue, thick]
      (3.8,0.8) -- (5.5,0.8);
    \node[below, codexblue, font=\tiny, align=center]
      at (4.6,0.7) {Rarefaction\\(low density)};

    % Oscillation direction
    \draw[<->, gray, thick]
      (10,2.2) -- (11.4,2.2);
    \node[right, gray, font=\tiny, align=center] at (11.5,2.2)
      {Oscillation\\direction};
  \end{scope}

\end{tikzpicture}
\end{center}
\end{diagbox}

% ============================================================
\section{Wave Terminology and Definitions}
% ============================================================

\begin{definitionbox}[Wave Parameters]
\begin{itemize}[noitemsep]
  \item \textbf{Amplitude} $A$: Maximum displacement
  of a particle from its equilibrium position. Unit: metre (m).

  \item \textbf{Wavelength} $\lambda$: Distance
  between two successive points that are in phase
  (e.g.\ crest to crest, trough to trough).
  Unit: metre (m).

  \item \textbf{Period} $T$: Time for one complete
  oscillation of a particle. Unit: second (s).

  \item \textbf{Frequency} $f$: Number of complete
  oscillations per second.
  $f = 1/T$. Unit: hertz (Hz).

  \item \textbf{Wave speed} $v$: Speed at which
  the wave pattern (disturbance) travels through
  the medium. Unit: m\,s$^{-1}$.

  \item \textbf{Phase}: Describes the stage of
  the oscillation cycle a particle is at (as an
  angle in degrees or radians).

  \item \textbf{Phase difference}: The fraction
  of a cycle by which two points on the wave are
  offset from each other.
\end{itemize}
\end{definitionbox}

\begin{diagbox}[Wave Profile: Key Parameters Labelled]
\begin{center}
\begin{tikzpicture}[font=\small]
  % ------------------------------------------------------------
  % Wave Profile: Key Parameters Labelled (computed, not eyeballed)
  % ------------------------------------------------------------
  % Conventions:
  %   y(x) = A\,\sin(kx) with k in rad per (abstract x-unit).
  % TikZ's sin() expects degrees, so we plot sin(deg(k*x)).
  %
  % This version fixes containment by *linking* the panel width and the plot
  % scale via explicit named quantities (\panelwidth, \wavelength, \xmargin).
  % The visible x-range is:  [crest1 - margin, crest2 + margin]
  % so the left/right padding is identical, and nothing can overflow.

  % --- Parameters (edit these; everything else follows) ---
  \pgfmathsetmacro{\A}{0.9}        % amplitude
  \pgfmathsetmacro{\k}{0.503}      % wavenumber (rad per x-unit)

  % --- Derived wave geometry (in abstract x-units) ---
  \pgfmathsetmacro{\wavelength}{2*pi/\k}          % \lambda
  \pgfmathsetmacro{\xCrestOne}{(pi/2)/\k}         % crest #1 at kx=pi/2
  \pgfmathsetmacro{\xCrestTwo}{\xCrestOne + \wavelength} % crest #2
  \pgfmathsetmacro{\xTrough}{(3*pi/2)/\k}         % trough at kx=3pi/2

  % --- Panel + padding (the "contract" that prevents overflow) ---
  \newlength{\panelwidth}
  \setlength{\panelwidth}{0.92\linewidth} % set to match your book/panel width

  \pgfmathsetmacro{\xmargin}{0.20*\wavelength}    % equal left/right x-padding
  \pgfmathsetmacro{\xMin}{\xCrestOne - \xmargin}
  \pgfmathsetmacro{\xMax}{\xCrestTwo + \xmargin}

  % Compute x-scale from the *fixed* panel width so exactly (\lambda+2m) fits.
  \pgfmathsetlengthmacro{\xunit}{\panelwidth/(\wavelength + 2*\xmargin)}

  % Vertical padding for labels/equations (in y-units; keep even top/bottom).
  % The wavelength label sits above the arrow; the previous bound clipped
  % its ascenders at the top of the surrounding diagbox.
  \pgfmathsetmacro{\yPadTop}{1.55}
  \pgfmathsetmacro{\yPadBottom}{2.05}

  % Annotation heights/offsets
  \pgfmathsetmacro{\yLam}{\A + 0.82}       % clear the crest label while retaining top padding
  \pgfmathsetmacro{\yEqText}{-\A - 1.45}   % equation block baseline (preserved)

  % --- Styles / font sizes ---
  \tikzset{
    pointlabel/.style={font=\small},
    paramlabel/.style={font=\small},
    eqntext/.style={font=\footnotesize, text=gray},
    amparrow/.style={<->, codexred, very thick},
    wave/.style={codexblue, very thick},
  }

  % Apply the computed x-scale (y stays in normal TikZ units).
  \begin{scope}[x=\xunit]

    % These coordinates must be defined inside the scaled scope. Defining
    % them outside made the second crest ignore \xunit and escape the panel.
    \coordinate (C1)  at (\xCrestOne, \A);
    \coordinate (C2)  at (\xCrestTwo, \A);
    \coordinate (T1)  at (\xTrough,  -\A);
    \coordinate (EqC1) at (\xCrestOne, 0);
    \coordinate (EqT1) at (\xTrough,    0);

    % Explicit bounding box = panel content + even padding (prevents overflow)
    \path[use as bounding box]
      (\xMin, -\A-\yPadBottom) rectangle (\xMax, \A+\yPadTop);

    % --- Wave curve (domain matches the bounding box x-range) ---
    \draw[wave, domain=\xMin:\xMax, samples=300]
      plot (\x, {\A*sin(deg(\k*\x))});

    % --- Equilibrium (mean) line (full curve width; preserved) ---
    \draw[gray, dashed] (\xMin,0) -- (\xMax,0);

    % --- Crest markers (exact extrema) ---
    \fill[codexgreen] (C1) circle (4pt);
    \fill[codexgreen] (C2) circle (4pt);

    % --- Labels: crest & trough (centered; preserved) ---
    \node[pointlabel, codexblue, above=3pt] at (C1) {Crest};
    \node[pointlabel, codexblue, below=3pt] at (T1) {Trough};

    % --- Amplitude arrows (preserved style/extent) ---
    \draw[amparrow] (EqC1) -- (C1);
    \node[paramlabel, codexred, right=4pt] at ($(EqC1)!0.5!(C1)$)
      {$A$ (amplitude)};

    \draw[amparrow] (EqT1) -- (T1);
    \node[paramlabel, codexred, right=4pt] at ($(EqT1)!0.5!(T1)$)
      {$-A$};

    % --- Wavelength arrow (fixed) ---
    % Double-headed (<->), arrowheads flush with crest marker x-coordinates,
    % and raised so it doesn't overlap the "Crest" label.
    \draw[<->, codexgold, very thick] (\xCrestOne, \yLam) -- (\xCrestTwo, \yLam);
    \draw[dashed, black!55] (C1) -- (\xCrestOne, \yLam);
    \draw[dashed, black!55] (C2) -- (\xCrestTwo, \yLam);
    \node[paramlabel, codexgold, above=2pt, font=\footnotesize,
      align=center, text width=0.82\panelwidth]
      at ($ (\xCrestOne,\yLam)!0.5!(\xCrestTwo,\yLam) $)
      {$\lambda$ (wavelength: crest to next crest)};

    % --- Defining equations (centered; preserved spacing) ---
    \node[eqntext] at ($ (\xMin,\yEqText)!0.5!(\xMax,\yEqText) $)
      {$T$ = time for one complete oscillation};
    \node[eqntext] at ($ (\xMin,\yEqText-0.42)!0.5!(\xMax,\yEqText-0.42) $)
      {$f = 1/T$, \quad $v = f\lambda$};

  \end{scope}
\end{tikzpicture}
\end{center}
\end{diagbox}

\begin{lawbox}[The Wave Equation]
The fundamental relationship between wave speed,
frequency and wavelength:
\[
  v = f\lambda = \frac{\lambda}{T}
\]
This applies to all waves --- mechanical waves,
electromagnetic waves, sound, light.
\end{lawbox}

\begin{examplebox}[17.1]
\stepcounter{workedexample}
A sound wave in air has frequency $440\ \text{Hz}$
(the musical note A) and travels at
$340\ \text{m\,s}^{-1}$.\\
(a) Calculate the wavelength.\\
(b) Calculate the period.\\
(c) Radio waves travel at $3.0\times10^8\ \text{m\,s}^{-1}$.
A radio station broadcasts at $99.3\ \text{MHz}$.
Calculate the wavelength of this broadcast.
\end{examplebox}

\begin{solutionbox}
\textbf{(a)} Wavelength of sound:
\[
  \lambda = \frac{v}{f}
  = \frac{340}{440}
  = 0.773\ \text{m} \approx 77.3\ \text{cm}
\]

\textbf{(b)} Period:
\[
  T = \frac{1}{f} = \frac{1}{440}
  = 2.27\times10^{-3}\ \text{s} = 2.27\ \text{ms}
\]

\textbf{(c)} Radio wave wavelength:
\[
  f = 99.3\ \text{MHz} = 99.3\times10^6\ \text{Hz}
\]
\[
  \lambda = \frac{v}{f}
  = \frac{3.0\times10^8}{99.3\times10^6}
  = 3.02\ \text{m}
\]
\end{solutionbox}

% ============================================================
\section{Properties of Waves}
% ============================================================

\subsection{Reflection}

When a wave meets a boundary between two media,
part of its energy is reflected back. The wave
bounces off the boundary.

\begin{lawbox}[Law of Reflection]
The angle of incidence equals the angle of
reflection, measured from the normal to the boundary:
\[
  \theta_i = \theta_r
\]
The incident wave, reflected wave and normal are all
in the same plane.
\end{lawbox}

\begin{diagbox}[Wave Reflection]
\begin{center}
\begin{tikzpicture}[scale=1.0, font=\small]

  % Reflecting surface
  \fill[gray!30] (0,-0.3) rectangle (8,0);
  \draw[very thick] (0,0) -- (8,0);
  \node[below, gray] at (4,-0.4) {Reflecting surface};

  % Normal
  \draw[dashed, gray, thick] (4,0) -- (4,3.5);
  \node[above, gray, font=\tiny] at (4,3.6) {Normal};

  % Incident wave ray
  \draw[->, very thick, codexblue, line width=2pt]
    (1.0,3.0) -- (4.0,0);
  \node[above left, codexblue] at (2.0,1.8)
    {Incident ray};

  % Reflected wave ray
  \draw[->, very thick, codexred, line width=2pt]
    (4.0,0) -- (7.0,3.0);
  \node[above right, codexred] at (6.0,1.8)
    {Reflected ray};

  % Angle of incidence
  \draw[codexblue, thick] (4,0.8) arc(90:147:0.8);
  \node[codexblue, font=\small] at (3.3,1.0)
    {$\theta_i$};

  % Angle of reflection
  \draw[codexred, thick] (4,0.8) arc(90:33:0.8);
  \node[codexred, font=\small] at (4.7,1.0)
    {$\theta_r$};

  % Law
  \node[font=\normalsize, codexgold]
    at (4,-1.0) {$\theta_i = \theta_r$};

\end{tikzpicture}
\end{center}
\end{diagbox}

\subsection{Refraction}

When a wave crosses from one medium into another
where the wave speed is different, the wavelength
changes, and the wave bends. The frequency remains
the same.

\begin{lawbox}[Snell's Law of Refraction]
\[
  \frac{\sin\theta_i}{\sin\theta_r}
  = \frac{v_1}{v_2}
  = \frac{\lambda_1}{\lambda_2}
  = n_{21}
\]
where $n_{21}$ is the refractive index of medium 2
relative to medium 1.

A wave slows down when entering a denser medium:
$v_2 < v_1 \implies \theta_r < \theta_i$
(bends toward the normal).
\end{lawbox}

\subsection{Diffraction}

When a wave passes through an opening or around an
obstacle, it spreads out into the region behind the
obstacle. This is \textbf{diffraction}.

\begin{keybox}[Diffraction is Most Significant When]
\[
  \lambda \approx d \quad \text{(wavelength
  $\approx$ gap/obstacle size)}
\]
If $\lambda \ll d$: little diffraction (wave travels
nearly straight). If $\lambda \approx d$ or
$\lambda > d$: significant spreading.
\end{keybox}

\begin{diagbox}[Diffraction Through a Gap]
\begin{center}
\begin{tikzpicture}[scale=0.68, font=\small]

  % === Wide gap (little diffraction) ===
  \begin{scope}[shift={(0,0)}]
    \node[font=\bfseries\small, align=center, text width=4.7cm] at (4,7.1)
      {Wide gap: little spreading};

    % Barrier
    \fill[gray!60] (3.5,-0.3) rectangle (4.5,2.2);
    \fill[gray!60] (3.5,3.5) rectangle (4.5,5.3);
    \draw[thick] (3.5,0) -- (3.5,2.2);
    \draw[thick] (3.5,5.3) -- (3.5,3.5);
    \draw[thick] (4.5,0) -- (4.5,2.2);
    \draw[thick] (4.5,5.3) -- (4.5,3.5);

    % Incoming straight waves
    \foreach \y in {0.5,1.0,1.5,2.0,2.5,3.0,3.5,4.0,4.5,5.0} {
      \draw[codexblue, thick] (0,\y) -- (3.5,\y);
    }

    % Transmitted waves (mostly straight, slight diffraction at edges)
    \foreach \y in {2.4,2.7,3.0,3.3} {
      \draw[codexblue, thick] (4.5,\y) -- (8,\y);
    }
    % Edge diffraction (small)
    \draw[codexblue!50, thick]
      (4.5,2.2) .. controls (5.5,2.2) and (6,1.8) .. (8,1.6);
    \draw[codexblue!50, thick]
      (4.5,3.5) .. controls (5.5,3.5) and (6,3.9) .. (8,4.1);

    % Gap label
    \draw[<->, gray] (4.7,2.2) -- (4.7,3.5);
    \node[right, gray, font=\tiny, fill=white, inner sep=1pt]
      at (4.8,2.85) {$d$};
  \end{scope}

  % === Narrow gap (strong diffraction) ===
  \begin{scope}[shift={(9.5,0)}]
    \node[font=\bfseries\small, align=center, text width=4.7cm] at (4,7.1)
      {Narrow gap: strong spreading};

    % Barrier
    \fill[gray!60] (3.5,-0.3) rectangle (4.5,2.5);
    \fill[gray!60] (3.5,3.1) rectangle (4.5,5.3);
    \draw[thick] (3.5,0) -- (3.5,2.5);
    \draw[thick] (3.5,5.3) -- (3.5,3.1);
    \draw[thick] (4.5,0) -- (4.5,2.5);
    \draw[thick] (4.5,5.3) -- (4.5,3.1);

    % Incoming straight waves
    \foreach \y in {0.5,1.0,1.5,2.0,2.5,3.0,3.5,4.0,4.5,5.0} {
      \draw[codexblue, thick] (0,\y) -- (3.5,\y);
    }

    % Circular/spreading waves after gap (diffraction)
    \foreach \r in {0.5,1.0,1.5,2.0,2.5,3.0,3.5} {
      \draw[codexblue, thick]
        (4.5,2.8) arc(-90:90:\r);
    }

    % Gap label
    \draw[<->, gray] (4.7,2.5) -- (4.7,3.1);
    \node[right, gray, font=\tiny, fill=white, inner sep=1pt]
      at (4.8,2.8) {$d$};
  \end{scope}

\end{tikzpicture}
\end{center}
\end{diagbox}

% ============================================================
\section{Superposition and Interference}
% ============================================================

\begin{lawbox}[Principle of Superposition]
When two or more waves meet at a point, the
resultant displacement at that point is the
\textbf{algebraic sum} of the individual
displacements due to each wave separately:
\[
  y_{total} = y_1 + y_2 + \cdots
\]
The waves pass through each other undisturbed ---
after overlapping, each continues as if the other
were not there.
\end{lawbox}

\begin{definitionbox}[Constructive and Destructive
Interference]
\textbf{Constructive interference:} Two waves
arrive \textit{in phase} (phase difference $= 0°$,
$360°$, $720°$, \ldots; path difference $= 0$,
$\lambda$, $2\lambda$, \ldots). Displacements add.
Resultant amplitude is maximum ($= A_1 + A_2$).

\vspace{0.4cm}

\textbf{Destructive interference:} Two waves arrive
\textit{in antiphase} (phase difference $= 180°$,
$540°$, \ldots; path difference $= \lambda/2$,
$3\lambda/2$, \ldots). Displacements subtract.
Resultant amplitude is minimum ($= |A_1 - A_2|$;
zero if $A_1 = A_2$).
\end{definitionbox}

\begin{diagbox}[Constructive and Destructive Interference]
\begin{center}
\begin{tikzpicture}[scale=0.85, font=\small]

  % === CONSTRUCTIVE ===
  \begin{scope}[shift={(0,6)}]
    \node[font=\bfseries, codexgreen] at (6,2.5)
      {Constructive Interference (in phase)};

    % Wave 1
    \draw[codexblue, thick, domain=0:12, samples=80]
      plot ({\x},
      {0.6*sin(deg(\x*0.52))+1.5});
    \node[left, codexblue, font=\tiny] at (0,1.5)
      {Wave 1};

    % Wave 2
    \draw[codexred, thick, domain=0:12, samples=80]
      plot ({\x},
      {0.6*sin(deg(\x*0.52))+0.5});
    \node[left, codexred, font=\tiny] at (0,0.5)
      {Wave 2};

    % Resultant
    \draw[codexgreen, very thick, domain=0:12, samples=80]
      plot ({\x},
      {1.2*sin(deg(\x*0.52))-0.8});
    \node[left, codexgreen, font=\small, align=center] at (0,-0.8)
      {Resultant\\($2A$)};
  \end{scope}

  % === DESTRUCTIVE ===
  \begin{scope}[shift={(0,0)}]
    \node[font=\bfseries, codexred] at (6,2.5)
      {Destructive Interference (antiphase)};

    % Wave 1
    \draw[codexblue, thick, domain=0:12, samples=80]
      plot ({\x},
      {0.6*sin(deg(\x*0.52))+1.5});
    \node[left, codexblue, font=\tiny] at (0,1.5)
      {Wave 1};

    % Wave 2 (antiphase: add pi to phase)
    \draw[codexred, thick, domain=0:12, samples=80]
      plot ({\x},
      {0.6*sin(deg(\x*0.52)+180)+0.5});
    \node[left, codexred, font=\tiny, align=center] at (0,0.5)
      {Wave 2\\(antiphase)};

    % Resultant (zero)
    \draw[codexgold, very thick, domain=0:12]
      plot ({\x},{-0.8});
    \node[left, codexgold, font=\small, align=center] at (0,-0.8)
      {Resultant\\$= 0$};
  \end{scope}

\end{tikzpicture}
\end{center}
\end{diagbox}

% ============================================================
\section{Stationary (Standing) Waves}
% ============================================================

When two waves of the same frequency, amplitude and
wavelength travel in opposite directions and overlap,
they produce a \textbf{stationary (standing) wave}.
The wave pattern appears to stand still --- some
points always have zero displacement
(\textbf{nodes}), while others oscillate with
maximum amplitude (\textbf{antinodes}).

\begin{processbox}[How a Standing Wave Forms]
Consider two waves travelling in opposite directions:
\[
  y_1 = A\sin(kx - \omega t)
  \quad \text{and} \quad
  y_2 = A\sin(kx + \omega t)
\]
Adding by superposition:
\[
  y = y_1 + y_2 = 2A\sin(kx)\cos(\omega t)
\]
This is NOT a travelling wave. The spatial part
$\sin(kx)$ and the temporal part $\cos(\omega t)$
are \textit{separate}. Every point oscillates with
the same frequency, but with amplitude that depends
on position: $A(x) = 2A\sin(kx)$.

\textbf{Nodes:} where $\sin(kx) = 0$:
$x = 0, \lambda/2, \lambda, 3\lambda/2, \ldots$

\textbf{Antinodes:} where $|\sin(kx)| = 1$:
$x = \lambda/4, 3\lambda/4, 5\lambda/4, \ldots$
\end{processbox}

\begin{diagbox}[Standing Wave: Nodes and Antinodes]
\begin{center}
\begin{tikzpicture}[scale=1.0, font=\small]

  % Standing wave pattern (multiple time snapshots)
  \draw[gray, dashed] (0,0) -- (12.5,0)
    node[right, font=\tiny]{Equilibrium};

  % Several snapshots of the standing wave
  \foreach \phase/\col/\opacity in {
    0/codexblue/100,
    30/codexblue/60,
    60/codexblue/30,
    90/gray/20} {
    \draw[\col!\opacity, thick, domain=0:12.56, samples=100]
      plot ({\x},
      {1.2*sin(deg(\x*0.5))*cos(deg(\phase))});
  }
  \draw[codexblue, very thick, domain=0:12.56, samples=100]
    plot ({\x},{1.2*sin(deg(\x*0.5))});

  % Nodes (marked with N)
  \foreach \x in {0, 6.28, 12.56} {
    \fill[codexred] (\x,0) circle (5pt);
    \node[below, codexred, font=\bfseries\small]
      at (\x,-0.3) {N};
  }

  % Antinodes (marked with A)
  \foreach \x/\y in {3.14/1.2, 9.42/-1.2} {
    \fill[codexgreen] (\x,\y) circle (5pt);
    \node[above, codexgreen, font=\bfseries\small]
      at (\x,{\y+0.15}) {A};
  }

  % Wavelength labels
  \draw[<->, codexgold, thick] (0,-1.65) -- (6.28,-1.65);
  \node[below, codexgold] at (3.14,-1.75)
    {$\lambda/2$};
  \draw[<->, codexgold, thick] (6.28,-1.65) -- (12.56,-1.65);
  \node[below, codexgold] at (9.42,-1.75)
    {$\lambda/2$};

  % Node-antinode distance
  \draw[<->, gray, thick] (0,-2.35) -- (3.14,-2.35);
  \node[below, gray, font=\tiny] at (1.57,-2.45)
    {$\lambda/4$};

  % Compact key, above the plotted extrema.
  \node[codexred, font=\tiny] at (6.28,2.25)
    {N = node (always zero displacement)};
  \node[codexgreen, font=\tiny] at (6.28,1.9)
    {A = antinode (maximum displacement)};

\end{tikzpicture}
\end{center}
\end{diagbox}

\begin{keybox}[Key Properties of Standing Waves]
\begin{itemize}[noitemsep]
  \item \textbf{Nodes:} points of zero displacement
  at all times. Separation between adjacent nodes
  $= \lambda/2$.
  \item \textbf{Antinodes:} points of maximum
  displacement. Separation between adjacent
  antinodes $= \lambda/2$.
  \item Node to adjacent antinode $= \lambda/4$.
  \item All particles between two adjacent nodes
  oscillate \textit{in phase} with each other.
  \item Particles on opposite sides of a node
  oscillate in \textit{antiphase}.
  \item \textbf{No energy is transferred} along
  the medium in a standing wave.
\end{itemize}
\end{keybox}

% ============================================================
\section{Standing Waves on Strings}
% ============================================================

When a string is fixed at both ends and vibrated,
standing waves form. The ends are \textbf{nodes}
(fixed points cannot move). The allowed wavelengths
are determined by the condition that an integer
number of half-wavelengths must fit between the
fixed ends.

\begin{lawbox}[Harmonics on a String Fixed at Both Ends]
For a string of length $L$ fixed at both ends,
the allowed wavelengths and frequencies are:

\[
  \lambda_n = \frac{2L}{n}, \quad
  f_n = \frac{nv}{2L} = nf_1
  \quad (n = 1, 2, 3, \ldots)
\]

where $v$ is the wave speed on the string,
$n = 1$ is the \textbf{fundamental} (first harmonic),
and $n = 2, 3, \ldots$ are the
\textbf{overtones} (harmonics).

The wave speed on a stretched string:
\[
  v = \sqrt{\frac{T}{\mu}}
\]
where $T$ is the tension (N) and
$\mu$ is the mass per unit length (kg\,m$^{-1}$).
\end{lawbox}

\begin{diagbox}[Harmonics on a String Fixed at Both Ends]
\begin{center}
\begin{tikzpicture}[scale=0.9, font=\small]

  % --- Fundamental (n=1) ---
  \begin{scope}[shift={(0,6)}]
    \node[left, font=\bfseries] at (0,0.5)
      {$n=1$\\Fund.};
    \draw[very thick] (0,0) -- (0,1);
    \draw[very thick] (8,0) -- (8,1);
    \draw[codexblue, very thick, domain=0:8, samples=60]
      plot ({\x},{0.7*sin(deg(\x*3.14/8))+0.15});
    \fill[codexred] (0,0.15) circle (4pt);
    \fill[codexred] (8,0.15) circle (4pt);
    \fill[codexgreen] (4,0.85) circle (4pt);
    \node[right] at (8.3,0.5)
      {$\lambda_1 = 2L$, $f_1 = v/2L$};
    \node[below, codexred, font=\tiny] at (0,0.05) {N};
    \node[below, codexred, font=\tiny] at (8,0.05) {N};
    \node[above, codexgreen, font=\tiny] at (4,0.85) {A};
  \end{scope}

  % --- 2nd harmonic (n=2) ---
  \begin{scope}[shift={(0,4)}]
    \node[left, font=\bfseries] at (0,0.5)
      {$n=2$\\2nd};
    \draw[very thick] (0,0) -- (0,1);
    \draw[very thick] (8,0) -- (8,1);
    \draw[codexred, very thick, domain=0:8, samples=80]
      plot ({\x},{0.7*sin(deg(\x*6.28/8))+0.15});
    \fill[codexred] (0,0.15) circle (4pt);
    \fill[codexred] (4,0.15) circle (4pt);
    \fill[codexred] (8,0.15) circle (4pt);
    \fill[codexgreen] (2,0.85) circle (4pt);
    \fill[codexgreen] (6,-0.55) circle (4pt);
    \node[right] at (8.3,0.5)
      {$\lambda_2 = L$, $f_2 = 2f_1$};
    \foreach \x in {0,4,8} {
      \node[below, codexred, font=\tiny] at (\x,0.05) {N};
    }
  \end{scope}

  % --- 3rd harmonic (n=3) ---
  \begin{scope}[shift={(0,2)}]
    \node[left, font=\bfseries] at (0,0.5)
      {$n=3$\\3rd};
    \draw[very thick] (0,0) -- (0,1);
    \draw[very thick] (8,0) -- (8,1);
    \draw[codexgreen, very thick, domain=0:8, samples=80]
      plot ({\x},{0.7*sin(deg(\x*9.42/8))+0.15});
    \fill[codexred] (0,0.15) circle (4pt);
    \fill[codexred] (2.67,0.15) circle (4pt);
    \fill[codexred] (5.33,0.15) circle (4pt);
    \fill[codexred] (8,0.15) circle (4pt);
    \fill[codexgreen] (1.33,0.85) circle (4pt);
    \fill[codexgreen] (4,-0.55) circle (4pt);
    \fill[codexgreen] (6.67,0.85) circle (4pt);
    \node[right] at (8.3,0.5)
      {$\lambda_3 = 2L/3$, $f_3 = 3f_1$};
  \end{scope}

  % Length label
  \draw[<->, gray, thick] (0,0.8) -- (8,0.8);
  \node[above, gray] at (4,0.9) {$L$};

  % Legend
  \node[codexred, font=\tiny, align=center] at (4,-0.45)
    {\ding{108} Node (N)};
  \node[codexgreen, font=\tiny, align=center] at (4,-0.8)
    {\ding{108} Antinode (A)};

\end{tikzpicture}
\end{center}
\end{diagbox}

% ============================================================
\section{Standing Waves in Pipes}
% ============================================================

Sound waves can form standing waves inside pipes
(organ pipes, flutes, the human vocal tract).
The boundary conditions depend on whether the pipe
is open or closed at each end.

\textbf{Open end:} displacement antinode
(pressure node) --- air is free to move.\\
\textbf{Closed end:} displacement node
(pressure antinode) --- air cannot move.

\begin{lawbox}[Standing Waves in Pipes]
\textbf{Open pipe} (open at both ends):
\[
  f_n = \frac{nv}{2L}, \quad n = 1, 2, 3, \ldots
\]
Both fundamental and all harmonics are present.

\vspace{0.4cm}

\textbf{Closed pipe} (closed at one end, open at other):
\[
  f_n = \frac{nv}{4L}, \quad n = 1, 3, 5, \ldots
  \text{ (odd harmonics only)}
\]
Only odd harmonics are present.
The fundamental frequency of a closed pipe of
length $L$ equals the fundamental of an open pipe
of length $2L$.
\end{lawbox}

\begin{diagbox}[Harmonics in Open and Closed Pipes]
\begin{center}
\begin{tikzpicture}[scale=0.8, font=\small]

  % --- OPEN PIPE ---
  \begin{scope}[shift={(0,0)}]
    \node[font=\bfseries, align=center] at (4,7.5)
      {Open Pipe};
    \node[font=\tiny, gray] at (4,7.15)
      {Air-displacement patterns (schematic)};

    % Pipe walls
    \draw[very thick, gray] (0,0) -- (0,7);
    \draw[very thick, gray] (8,0) -- (8,7);

    % Harmonics (n = 1,2,3). Both open ends are displacement antinodes.
    \foreach \n/\col/\ybase in {1/codexblue/6.0, 2/codexred/4.0, 3/codexgreen/2.0} {
      \draw[\col, very thick, domain=0:8, samples=80]
        plot ({\x},{0.6*cos(deg(\x*\n*3.14159/8))+\ybase});
      \node[left, font=\tiny, text=\col] at (-0.2,\ybase)
        {$n=\n$};
    }

    % n=1: antinodes at both open ends, node at the centre.
    \fill[codexgreen] (0,6.6) circle (4pt);
    \fill[codexgreen] (8,5.4) circle (4pt);
    \fill[codexred] (4,6.0) circle (4pt);
    % n=2: antinodes at both ends and centre; nodes at quarter points.
    \fill[codexgreen] (0,4.6) circle (4pt);
    \fill[codexgreen] (4,3.4) circle (4pt);
    \fill[codexgreen] (8,4.6) circle (4pt);
    \foreach \x in {2,6} {\fill[codexred] (\x,4.0) circle (4pt);}
    % n=3: four antinodes and three interior nodes.
    \fill[codexgreen] (0,2.6) circle (4pt);
    \fill[codexgreen] (2.67,1.4) circle (4pt);
    \fill[codexgreen] (5.33,2.6) circle (4pt);
    \fill[codexgreen] (8,1.4) circle (4pt);
    \foreach \x in {1.33,4,6.67} {\fill[codexred] (\x,2.0) circle (4pt);}

    % Labels
    \node[below, font=\small] at (4,-0.3)
      {$L$};
    \draw[<->, gray] (0,-0.5) -- (8,-0.5);
  \end{scope}

  % --- CLOSED PIPE ---
  \begin{scope}[shift={(10,0)}]
    \node[font=\bfseries, align=center] at (4,7.5)
      {Closed Pipe};

    % Pipe walls
    \draw[very thick, gray] (0,0) -- (0,7);
    \draw[very thick, gray] (8,0) -- (8,7);
    % Closed end (left)
    \fill[gray!60] (-0.3,0) rectangle (0,7);
    \draw[very thick] (-0.3,0) rectangle (0,7);
    \node[below, font=\tiny] at (0,-1.0) {Closed};
    \node[below, font=\tiny] at (8,-1.0) {Open};

    % n=1 (fundamental: quarter wave)
    \draw[codexblue, very thick, domain=0:8, samples=80]
      plot ({\x},{0.7*sin(deg(\x*1.57/8))+6.0});
    \fill[codexred] (0,6.0) circle (4pt);
    \fill[codexgreen] (8,6.7) circle (4pt);
    \node[left, font=\tiny, codexblue] at (-0.5,6.0)
      {$n=1$};
    \node[font=\tiny, align=center] at (6.1,6.9)
      {$f_1 = v/4L$};

    % n=3 (3rd harmonic: 3/4 wave)
    \draw[codexred, very thick, domain=0:8, samples=80]
      plot ({\x},{0.7*sin(deg(\x*4.71/8))+4.0});
    \fill[codexred] (0,4.0) circle (4pt);
    \fill[codexgreen] (2.67,4.7) circle (3pt);
    \fill[codexred] (5.33,4.0) circle (3pt);
    \fill[codexgreen] (8,3.3) circle (4pt);
    \node[left, font=\tiny, codexred] at (-0.5,4.0)
      {$n=3$};
    \node[font=\tiny, align=center] at (6.1,4.9)
      {$f_3 = 3f_1$};

    % n=5
    \draw[codexgreen, very thick, domain=0:8, samples=80]
      plot ({\x},{0.7*sin(deg(\x*7.85/8))+2.0});
    \fill[codexred] (0,2.0) circle (4pt);
    \foreach \x in {3.2,6.4} {\fill[codexred] (\x,2.0) circle (3pt);}
    \fill[codexgreen] (1.6,2.7) circle (3pt);
    \fill[codexgreen] (4.8,1.3) circle (3pt);
    \fill[codexgreen] (8,2.7) circle (4pt);
    \node[left, font=\tiny, codexgreen] at (-0.5,2.0)
      {$n=5$};
    \node[font=\tiny, align=center] at (6.1,2.9)
      {$f_5 = 5f_1$\\only odd};

    \draw[<->, gray] (0,-0.5) -- (8,-0.5);
    \node[below, font=\small] at (4,-0.3) {$L$};
  \end{scope}

\end{tikzpicture}
\end{center}
\end{diagbox}

\begin{examplebox}[17.2]
\stepcounter{workedexample}
A guitar string of length $65\ \text{cm}$ and mass
per unit length $4.0\times10^{-3}\ \text{kg\,m}^{-1}$
is under tension $120\ \text{N}$.\\
(a) Calculate the wave speed on the string.\\
(b) Calculate the fundamental frequency.\\
(c) Calculate the frequency of the 3rd harmonic.\\
(d) What length would the string need to be to
produce a fundamental of $440\ \text{Hz}$?
\end{examplebox}

\begin{solutionbox}
\textbf{(a)} Wave speed:
\[
  v = \sqrt{\frac{T}{\mu}}
  = \sqrt{\frac{120}{4.0\times10^{-3}}}
  = \sqrt{30\,000}
  = 173.2\ \text{m\,s}^{-1}
\]

\textbf{(b)} Fundamental ($L = 0.65\ \text{m}$):
\[
  f_1 = \frac{v}{2L}
  = \frac{173.2}{2 \times 0.65}
  = \frac{173.2}{1.30}
  = 133.2\ \text{Hz}
\]

\textbf{(c)} 3rd harmonic:
\[
  f_3 = 3f_1 = 3 \times 133.2 = 399.7\ \text{Hz}
  \approx 400\ \text{Hz}
\]

\textbf{(d)} For $f_1 = 440\ \text{Hz}$:
\[
  L = \frac{v}{2f_1}
  = \frac{173.2}{2 \times 440}
  = \frac{173.2}{880}
  = 0.197\ \text{m} = 19.7\ \text{cm}
\]
The guitarist must shorten the string to $\approx 20\ \text{cm}$
by pressing a fret.
\end{solutionbox}

\begin{examplebox}[17.3]
\stepcounter{workedexample}
An organ pipe, open at both ends, has length
$0.680\ \text{m}$. Speed of sound $= 340\ \text{m\,s}^{-1}$.\\
(a) Find the fundamental frequency.\\
(b) List the next three harmonics.\\
(c) If the pipe is closed at one end (keeping
the same length), what is the new fundamental
frequency?
\end{examplebox}

\begin{solutionbox}
\textbf{(a)} Open pipe fundamental:
\[
  f_1 = \frac{v}{2L}
  = \frac{340}{2 \times 0.680}
  = \frac{340}{1.360}
  = 250\ \text{Hz}
\]

\textbf{(b)} Next three harmonics (open pipe has
all harmonics):
\begin{align*}
  f_2 &= 2 \times 250 = 500\ \text{Hz}\\
  f_3 &= 3 \times 250 = 750\ \text{Hz}\\
  f_4 &= 4 \times 250 = 1000\ \text{Hz}
\end{align*}

\textbf{(c)} Closed pipe fundamental (same length):
\[
  f_1' = \frac{v}{4L}
  = \frac{340}{4 \times 0.680}
  = \frac{340}{2.720}
  = 125\ \text{Hz}
\]
The closed pipe produces half the fundamental
frequency --- one octave lower --- and only odd
harmonics ($125$, $375$, $625$, $875\ \text{Hz}$, \ldots).
\end{solutionbox}

% ============================================================
\section{Electromagnetic Waves}
% ============================================================

\begin{definitionbox}[Electromagnetic Waves]
\textbf{Electromagnetic (EM) waves} are transverse
waves consisting of oscillating electric and magnetic
fields, perpendicular to each other and to the
direction of propagation. They:
\begin{itemize}[noitemsep]
  \item Require no medium (travel through vacuum)
  \item All travel at speed $c = 3.0\times10^8\ \text{m\,s}^{-1}$
  in vacuum
  \item Carry energy but no mass
  \item Obey $c = f\lambda$
  \item Can be polarised (transverse waves)
\end{itemize}
\end{definitionbox}

\begin{table}[H]
\centering
\caption{The Electromagnetic Spectrum}
\label{tab:em_spectrum}
\begin{tabular}{lll}
\toprule
\textbf{Type} & \textbf{Wavelength} & \textbf{Applications}\\
\midrule
Radio waves & $>1\ \text{mm}$ &
  Broadcasting, radar, WiFi\\
Microwaves & $1\ \text{mm}$--$1\ \text{m}$ &
  Microwave ovens, mobile phones\\
Infrared & $700\ \text{nm}$--$1\ \text{mm}$ &
  Remote controls, heat cameras\\
Visible light & $400$--$700\ \text{nm}$ &
  Vision, photography\\
Ultraviolet & $10$--$400\ \text{nm}$ &
  Sterilisation, vitamin D\\
X-rays & $0.01$--$10\ \text{nm}$ &
  Medical imaging, security\\
Gamma rays & $<0.01\ \text{nm}$ &
  Cancer treatment, sterilisation\\
\bottomrule
\end{tabular}
\end{table}

% ============================================================
\section{Chapter Summary}
% ============================================================

\begin{summarybox}
\textbf{Wave:} propagating disturbance transferring
energy without transferring matter.

\textbf{Transverse:} oscillation $\perp$ propagation
(light, string waves, water).
\textbf{Longitudinal:} oscillation $\parallel$
propagation (sound).

\textbf{Wave equation:} $v = f\lambda = \lambda/T$

\textbf{Wave properties:}
Reflection ($\theta_i = \theta_r$);
Refraction (bends at boundary, $v$ and $\lambda$ change,
$f$ constant);
Diffraction (spreads, most for $\lambda \approx d$);
Interference (constructive: path diff. $= n\lambda$;
destructive: path diff. $= (n+\tfrac{1}{2})\lambda$)

\textbf{Superposition:} $y = y_1 + y_2$

\textbf{Standing waves:}
Nodes (zero displacement) separated by $\lambda/2$.
Antinodes (max displacement) between nodes.

\textbf{String (fixed both ends):}
$f_n = nv/(2L)$, all harmonics.

\textbf{Open pipe:} $f_n = nv/(2L)$, all harmonics.

\textbf{Closed pipe:} $f_n = nv/(4L)$,
odd harmonics only.

\textbf{Wave speed on string:}
$v = \sqrt{T/\mu}$

\textbf{EM waves:} transverse, speed $c$
in vacuum, no medium needed.
\end{summarybox}

% ============================================================
\newpage
\begin{exercisebox}[17]

\subsection*{Section A: Multiple Choice}

\textbf{A1.} Which of the following is a
longitudinal wave?

\mcoption{A}{Light}
\mcoption{B}{Water ripples}
\mcoption{C}{Sound}
\mcoption{D}{Waves on a guitar string}
\ansref{Ch17-A1}

\vspace{0.4cm}

\textbf{A2.} A wave has speed $320\ \text{m\,s}^{-1}$
and wavelength $0.8\ \text{m}$. Its frequency is:

\mcoption{A}{$25.6\ \text{Hz}$}
\mcoption{B}{$256\ \text{Hz}$}
\mcoption{C}{$400\ \text{Hz}$}
\mcoption{D}{$40\ \text{Hz}$}
\ansref{Ch17-A2}

\vspace{0.4cm}

\textbf{A3.} Diffraction of a wave is most
significant when:

\mcoption{A}{The wavelength is much less than
the gap width}
\mcoption{B}{The wavelength is much greater than
the gap width}
\mcoption{C}{The wavelength is approximately
equal to the gap width}
\mcoption{D}{The wave speed is very high}
\ansref{Ch17-A3}

\vspace{0.4cm}

\textbf{A4.} Two waves of equal amplitude arrive
at a point with a path difference of $3\lambda/2$.
The interference is:

\mcoption{A}{Constructive}
\mcoption{B}{Destructive}
\mcoption{C}{Partial}
\mcoption{D}{No interference occurs}
\ansref{Ch17-A4}

\vspace{0.4cm}

\textbf{A5.} In a standing wave on a string, the
distance between a node and the nearest antinode is:

\mcoption{A}{$\lambda$}
\mcoption{B}{$\lambda/2$}
\mcoption{C}{$\lambda/4$}
\mcoption{D}{$2\lambda$}
\ansref{Ch17-A5}

\vspace{0.4cm}

\textbf{A6.} A closed pipe of length $L$ has
fundamental frequency $f$. An open pipe of the
same length has fundamental frequency:

\mcoption{A}{$f/2$}
\mcoption{B}{$f$}
\mcoption{C}{$2f$}
\mcoption{D}{$4f$}
\ansref{Ch17-A6}

\vspace{0.4cm}

\textbf{A7.} The speed of a wave on a stretched
string is doubled by:

\mcoption{A}{Doubling the length of the string}
\mcoption{B}{Halving the mass per unit length}
\mcoption{C}{Quadrupling the tension}
\mcoption{D}{Doubling the frequency}
\ansref{Ch17-A7}

\vspace{0.4cm}

\textbf{A8.} When a wave passes from a less dense
medium to a denser medium, which quantity remains
unchanged?

\mcoption{A}{Speed}
\mcoption{B}{Wavelength}
\mcoption{C}{Frequency}
\mcoption{D}{Amplitude}
\ansref{Ch17-A8}

\vspace{0.4cm}

\textbf{A9.} A standing wave is formed in a string
$1.5\ \text{m}$ long. If the wave speed on the
string is $300\ \text{m\,s}^{-1}$, the third
harmonic frequency is:

\mcoption{A}{$100\ \text{Hz}$}
\mcoption{B}{$200\ \text{Hz}$}
\mcoption{C}{$300\ \text{Hz}$}
\mcoption{D}{$600\ \text{Hz}$}
\ansref{Ch17-A9}

\vspace{0.4cm}

\textbf{A10.} Electromagnetic waves differ from
mechanical waves because electromagnetic waves:

\mcoption{A}{Travel faster than mechanical waves}
\mcoption{B}{Do not require a medium and can
travel through vacuum}
\mcoption{C}{Have shorter wavelengths}
\mcoption{D}{Cannot be reflected}
\ansref{Ch17-A10}

\vspace{0.6cm}

\subsection*{Section B: Theory and Calculations}

\textbf{B1.} (a) Explain the difference between
a transverse and a longitudinal wave, giving two
examples of each.\\[4pt]
(b) A water wave has wavelength $3.0\ \text{m}$
and period $2.5\ \text{s}$.\\
(i) Calculate the wave speed.\\
(ii) A boat 150 m from the source of the waves
will feel the waves after how many seconds?\\
(iii) The wave passes a fixed buoy. How many
complete waves pass the buoy in 30 s?
\hfill\ansref{Ch17-B1}

\vspace{0.4cm}

\textbf{B2.} A string of length $0.90\ \text{m}$
and mass $6.0\times10^{-3}\ \text{kg}$ is stretched
under a tension of $150\ \text{N}$.\\[4pt]
(a) Calculate the mass per unit length.\\
(b) Calculate the wave speed on the string.\\
(c) Calculate the fundamental frequency and the
frequencies of the 2nd, 3rd, and 4th harmonics.\\
(d) Sketch the standing wave pattern for the
3rd harmonic, labelling all nodes and antinodes.\\
(e) The tension is increased. Explain qualitatively
what happens to the fundamental frequency.
\hfill\ansref{Ch17-B2}

\vspace{0.4cm}

\textbf{B3.} A vibrating tuning fork of frequency
$512\ \text{Hz}$ is held above a closed pipe
(open at the top). The speed of sound is
$340\ \text{m\,s}^{-1}$.\\[4pt]
(a) Calculate the wavelength of the sound.\\
(b) Calculate the shortest length of pipe for
which resonance occurs (fundamental).\\
(c) Calculate the next two lengths at which
resonance occurs.\\
(d) If the same pipe were open at both ends,
what length would give resonance at the same
frequency (fundamental)?\\
(e) Explain why only odd harmonics are present
in a pipe closed at one end.
\hfill\ansref{Ch17-B3}

\vspace{0.4cm}

\textbf{B4.} Two speakers face each other, $3.0\ \text{m}$
apart, each emitting sound of frequency
$680\ \text{Hz}$ in phase. Speed of sound
$= 340\ \text{m\,s}^{-1}$.\\[4pt]
(a) Explain why a standing wave is set up between
the speakers.\\
(b) Calculate the wavelength.\\
(c) Calculate the positions of the nodes between
the speakers (measuring from one speaker).\\
(d) At what positions are the antinodes?\\
(e) A listener walks between the speakers along the
line joining them. Describe what they hear.
\hfill\ansref{Ch17-B4}

\vspace{0.4cm}

\textbf{B5.} (a) State the Principle of Superposition
and explain how it leads to constructive and
destructive interference.\\[4pt]
(b) Two coherent point sources $S_1$ and $S_2$ emit
waves of wavelength $4.0\ \text{cm}$ in phase.
Point P is $12.0\ \text{cm}$ from $S_1$ and
$14.0\ \text{cm}$ from $S_2$.\\
(i) Calculate the path difference at P.\\
(ii) State whether constructive or destructive
interference occurs at P.\\[4pt]
(c) Point Q is $12.0\ \text{cm}$ from $S_1$ and
$16.0\ \text{cm}$ from $S_2$.
Determine the type of interference at Q.\\[4pt]
(d) Explain what happens to the interference
pattern if the two sources are moved closer
together, keeping the wavelength constant.
\hfill\ansref{Ch17-B5}

\end{exercisebox}
